% AulaTeX - maqueta inicial de construcción descendente
% Contrato futuro: el Agente puede investigar, redactar, evaluar y compilar a partir de este archivo.
\documentclass[12pt]{article}
\usepackage[utf8]{inputenc}
\usepackage[T1]{fontenc}
\usepackage[spanish]{babel}
\usepackage{enumitem}
\usepackage{hyperref}
\usepackage{longtable}

\title{derecho-a-la-seguridad-social-lde}
\author{AulaTeX}
\date{\today}

\begin{document}
\maketitle

\noindent\textbf{Nodo editorial:} UnADM/derecho-a-la-seguridad-social-lde\\
\textbf{Padre editorial:} UnADM\\
\textbf{Modo de generación:} reforzar\\
\textbf{Destino:} UnADM/licenciatura-en-derecho-unadm/derecho-a-la-seguridad-social-lde\\
\textbf{Contrato futuro:} ready-for-agent

\section*{Ingesta base}
\begin{itemize}[leftmargin=*]
  \item Texto libre proporcionado: no
  \item Documento de apoyo proporcionado: no
\end{itemize}

\section*{Objetivo}
\begin{itemize}[leftmargin=*]
  \item Analizar el derecho a la seguridad social desde su fundamento constitucional, legal e institucional.
  \item Aplicar criterios jurídicos para resolver problemas de acceso y protección social en casos concretos.
  \item Completar la maqueta inicial de derecho-a-la-seguridad-social-lde.
  \item Analizar el derecho a la seguridad social desde su base constitucional, legal e institucional.
  \item Resolver problemas jurídicos de cobertura y protección social con argumentación sustentada.
  \item Resolver problemas juridicos de cobertura y proteccion social con argumentacion sustentada.
\end{itemize}

\section*{Competencias}
\begin{itemize}[leftmargin=*]
  \item Interpretación de normas de seguridad social.
  \item Argumentación jurídica con base en fuentes primarias y secundarias.
  \item Elaboración de conclusiones aplicables al ejercicio profesional.
  \item Interpretación sistemática de normas de seguridad social.
  \item Construcción de argumentos jurídicos con fuentes primarias y secundarias.
  \item Formulación de conclusiones aplicables a la práctica profesional.
  \item Interpretacion sistematica de normas de seguridad social.
  \item Construccion de argumentos juridicos con fuentes primarias y secundarias.
  \item Formulacion de conclusiones aplicables a la practica profesional.
\end{itemize}

\section*{Resultados esperados}
\begin{itemize}[leftmargin=*]
  \item Identifica marco normativo esencial y su jerarquía.
  \item Explica obligaciones del Estado e instituciones de seguridad social.
  \item Sostiene una postura jurídica propia con evidencia.
  \item Delimita correctamente el problema jurídico planteado.
  \item Identifica y jerarquiza fuentes normativas aplicables.
  \item Sostiene una postura jurídica propia con soporte verificable.
  \item Delimita correctamente el problema juridico planteado.
  \item Sostiene una postura juridica propia con soporte verificable.
  \item Integra conclusiones transferibles al ejercicio profesional.
\end{itemize}

\section*{Estructura sugerida}
\begin{itemize}[leftmargin=*]
  \item Introducción y planteamiento del problema.
  \item Marco conceptual y normativo.
  \item Análisis jurídico del caso o tema.
  \item Conclusiones y recomendaciones.
  \item Referencias bibliográficas.
  \item Datos de actividad y objetivo.
  \item Planteamiento del problema jurídico.
  \item Marco normativo y conceptual.
  \item Análisis y argumentación.
  \item Conclusiones.
  \item Referencias.
  \item Planteamiento del problema juridico.
  \item Analisis y argumentacion juridica.
\end{itemize}

\section*{Criterios de evaluación}
\begin{itemize}[leftmargin=*]
  \item Dominio conceptual y normativo.
  \item Rigor argumentativo.
  \item Pertinencia de fuentes.
  \item Claridad de redacción jurídica.
  \item Cumplimiento formal del formato LaTeX institucional.
  \item Dominio conceptual-normativo.
  \item Pertinencia y trazabilidad de fuentes.
  \item Claridad expositiva y técnica jurídica.
  \item Cumplimiento formal LaTeX institucional.
  \item Rigor argumentativo y tecnica juridica.
  \item Claridad expositiva.
\end{itemize}

\section*{Bibliografía requerida}
\begin{itemize}[leftmargin=*]
  \item Fuentes oficiales mexicanas vigentes en seguridad social.
  \item Programa analítico local de la materia.
  \item Bibliografía institucional UnADM pertinente al campo jurídico.
  \item Material institucional UnADM pertinente.
  \item Doctrina y jurisprudencia verificables según tema.
  \item Programa analitico local de la materia.
  \item Doctrina y jurisprudencia verificables segun tema.
  \item CPEUM, Ley del Seguro Social y Ley del ISSSTE en version oficial vigente.
\end{itemize}

\section*{Marcadores de investigación}
\begin{itemize}[leftmargin=*]
  \item Artículo 123 constitucional y seguridad social.
  \item Prestaciones, cobertura y sujetos protegidos en LSS e ISSSTE.
  \item Mecanismos de tutela y medios de defensa.
  \item Tendencias y desafíos contemporáneos del sistema.
  \item Definir fuentes, citas y vacíos de contenido antes de redactar.
  \item Artículo 123 constitucional y contenido prestacional.
  \item Régimen IMSS/ISSSTE: sujetos, prestaciones y contingencias.
  \item Medios de defensa en materia de seguridad social.
  \item Criterios jurisdiccionales recientes sobre acceso y suficiencia.
  \item Tensiones entre universalidad normativa y cobertura efectiva.
  \item Articulo 123 constitucional y contenido prestacional.
  \item Regimen IMSS/ISSSTE: sujetos, prestaciones y contingencias.
  \item Tension entre universalidad normativa y cobertura efectiva.
\end{itemize}

\section*{Indicaciones editoriales por entregable}

\subsection*{Plantilla}
\begin{itemize}[leftmargin=*]
  \item \\textbackslash\\{\\}section\\{Datos de la actividad\\}
  \item \\textbackslash\\{\\}section\\{Objetivo\\}
  \item \\textbackslash\\{\\}section\\{Problema juridico\\}
  \item \\textbackslash\\{\\}section\\{Marco normativo y conceptual\\}
  \item \\textbackslash\\{\\}section\\{Analisis y argumentacion\\}
  \item \\textbackslash\\{\\}section\\{Conclusiones\\}
  \item \\textbackslash\\{\\}bibliography\\{derecho-a-la-seguridad-social\\}
  \item Definir plantilla base con reglas editoriales, assets y referencias de estilo antes de redactar.
  \item \\textbackslash\\{\\}section\\{Planteamiento del problema juridico\\}
  \item \\textbackslash\\{\\}section\\{Analisis y argumentacion juridica\\}
  \item Prepara una materia con estructura apta para actividades, bibliografía y entregables académicos reutilizables.
  \item Definir una maqueta compatible con latexmk y rutas canónicas de AulaTeX.
  \item Mantener \\textbackslash\\{\\}bibliography apuntando al .bib local de la materia.
  \item Evitar paquetes innecesarios en esta fase; solo placeholders estructurales.
  \item Organizar cada entrega en: problema, marco conceptual-normativo, análisis, conclusión jurídica y referencias.
  \item Usar archivos canónicos de materia: reporte-*.tex, presentacion-*.tex y derecho-a-la-seguridad-social.bib.
  \item Separar contenido de investigación (fuentes) y redacción final (entregable).
  \item Trabaja con la memoria fundacional consolidada como fuente base.
\end{itemize}

\subsection*{Actividad}
\begin{itemize}[leftmargin=*]
  \item Definir consigna y alcance de unidad.
  \item Incluir instrucciones de evidencia y formato de entrega.
  \item Reservar bloque para postura juridica personal sustentada.
  \item Desarrollar la actividad respetando objetivo, criterios de evaluación y marcadores de investigación.
  \item Definir consigna puntual, alcance y evidencia esperada por unidad.
  \item Incluir instrucciones de formato, extensión y criterio de evaluación.
  \item Reservar subsección para postura jurídica personal sustentada.
  \item Vincular explícitamente con marcadores de investigación.
  \item Incluir instrucciones de formato, extension y criterio de evaluacion.
  \item Reservar subseccion para postura juridica personal sustentada.
  \item Vincular explicitamente con marcadores de investigacion.
  \item Analizar el derecho a la seguridad social desde su fundamento constitucional, legal e institucional.
  \item Aplicar criterios jurídicos para resolver problemas de acceso y protección social en casos concretos.
  \item Dominio conceptual y normativo.
  \item Rigor argumentativo.
  \item Pertinencia de fuentes.
  \item Marco constitucional del derecho a la seguridad social en México.
  \item Leyes e instituciones clave (IMSS, ISSSTE y régimen aplicable).
  \item Principios de universalidad, progresividad, igualdad y no discriminación.
  \item Trabaja con la memoria fundacional consolidada como fuente base.
\end{itemize}

\subsection*{Reporte}
\begin{itemize}[leftmargin=*]
  \item Estructura IMRyC adaptada a ciencias juridicas.
  \item Agregar tabla breve de normas/fuentes aplicadas.
  \item Cerrar con implicaciones practicas del analisis.
  \item Estructurar el reporte con bibliografía requerida y cierre verificable.
  \item Usar secuencia: problema, fundamento, análisis, conclusión, referencias.
  \item Agregar tabla breve de normas/criterios aplicados.
  \item Cerrar con implicaciones prácticas del análisis jurídico.
  \item Verificar correspondencia total entre citas y .bib.
  \item Usar secuencia: problema, fundamento, analisis, conclusion, referencias.
  \item Cerrar con implicaciones practicas del analisis juridico.
  \item Portada y datos institucionales.
  \item Objetivo de actividad o unidad.
  \item Desarrollo: problema, fundamentos, análisis y postura.
  \item Redacción clara, precisa y sin relleno.
  \item Priorizar párrafos breves con conectores lógicos y terminología jurídica correcta.
  \item Evitar afirmaciones sin fuente o sin criterio jurídico explícito.
  \item No inventar referencias bibliográficas.
  \item Usar primero fuentes institucionales y normativas aplicables.
  \item Completar derecho-a-la-seguridad-social.bib con entradas verificables en fase de investigación.
  \item JSON y estructura editorial válidos antes de bajar a fase de agente.
  \item Coherencia entre objetivo, competencias, criterios y producto solicitado.
  \item Trazabilidad mínima de fuentes (normas, doctrina, material institucional).
\end{itemize}

\subsection*{Presentación}
\begin{itemize}[leftmargin=*]
  \item Diapositivas: contexto, problema, marco legal, hallazgos, cierre.
  \item Reducir texto y priorizar mapas normativos/sintesis.
  \item Incluir diapositiva final de referencias clave.
  \item Preparar la presentación con síntesis visual, continuidad editorial y evidencias clave.
  \item Estructura de diapositivas: contexto, problema, marco legal, hallazgos, cierre.
  \item Reducir texto; priorizar mapas normativos y síntesis visual.
  \item Incluir diapositiva final de referencias clave verificables.
  \item Mantener continuidad argumentativa con el reporte.
  \item Reducir texto y priorizar mapas normativos y sintesis visual.
  \item Identifica marco normativo esencial y su jerarquía.
  \item Explica obligaciones del Estado e instituciones de seguridad social.
  \item Sostiene una postura jurídica propia con evidencia.
  \item Redacción clara, precisa y sin relleno.
  \item Priorizar párrafos breves con conectores lógicos y terminología jurídica correcta.
  \item Incluye memoria, plan y maqueta inicial de materia.
  \item No incluye investigación exhaustiva ni redacción completa de actividades.
  \item Sintetiza visualmente la memoria editorial sin perder continuidad con plantilla, actividad y reporte.
\end{itemize}

\end{document}
